\begin{tikzpicture}[scale=1]

% === Parâmetros ===
\def\r{3}   % raio do círculo
\def\a{2}   % retas verticais y=±a

\pgfmathsetmacro{\zc}{sqrt(\r^2-\a^2)}     % altura z da interseção: sqrt(5)
\pgfmathsetmacro{\ang}{acos(\a/\r)}        % ângulo da interseção com y=a

  \begin{axis}[      
      xlabel={$y$}, ylabel={$z$},
      axis equal image,        % mantém escala y=z para o círculo não ficar ovalado
      xmin=-1.3*\r, xmax=1.3*\r, ymin=-1.3*\r, ymax=1.3*\r,
      xtick={-\r,-\a,\a,\r}, ytick={-\zc,\zc},
      yticklabels={$-\sqrt{5}$,$\sqrt{5}$},
      x tick label style={anchor=north east}, % rótulos abaixo e a esquerda do eixo horizontal
      y tick label style={anchor=south east}, % rótulos acima e a esquerda do eixo vertical
      every axis x label/.style={at={(ticklabel* cs:0.97)}, anchor=north},
      every axis y label/.style={at={(ticklabel* cs:0.97)}, anchor=east},
    ]

    % --- Região sombreada: dentro do círculo e à direita de y=a ---
    \path[fill=gray!30, opacity=0.6]
      (axis cs:\a,-\zc)     % Ponto inicial, na interseção inferior
      arc[start angle=-\ang, end angle=\ang, radius=\r] % Arco do círculo
      -- cycle;             % Fecha com o segmento da reta y=a

    % --- Linhas Tracejadas ---
    \addplot[Tracejado] coordinates {(0,\zc) (\a,\zc)};
    \addplot[Tracejado] coordinates {(0,-\zc) (\a,-\zc)};

    % --- Curvas ---
    \draw[thick, blue] (axis cs:0,0) circle [radius=\r];    % Círculo de raio r
    \addplot[thick, red] coordinates {(\a,-1.2*\r) (\a,1.2*\r)} node[pos=1, right, red] {$y=2$};
    \addplot[thick, red] coordinates {(-\a,-1.2*\r) (-\a,1.2*\r)} node[pos=1, left, red] {$y=-2$};

    % --- Pontos de interseção ---
    \addplot[only marks, black!60, mark=*] coordinates {(\a,\zc) (\a,-\zc) (-\a,\zc) (-\a,-\zc)};

  \end{axis}
  
\end{tikzpicture}
